% ============================================================================
% Evaluation Protocol Flaws in Facial Kinship Verification
%
% All numeric claims in this manuscript trace to results/honest/*.json and are
% reproducible via scripts/evaluation/run_kfold.py. Do not add a figure or a
% baseline row without a verifiable source.
% ============================================================================

\documentclass[journal,12pt]{IEEEtran}

\usepackage{amsmath,amssymb,amsfonts}
\usepackage{graphicx}
\usepackage{booktabs}
\usepackage{multirow}
\usepackage{url}
\usepackage[hidelinks]{hyperref}

\begin{document}

\title{Identity Leakage, Dataset Shortcuts, and a Negative Result for
Quantum-Inspired Metrics in Facial Kinship Verification}

\author{Sasi~Vaibhav% <-- add co-authors and affiliation before submission
\thanks{Code and data-processing scripts:
\url{https://github.com/svaibhav-7/Quantum_kinship_verification}}}

\maketitle

\begin{abstract}
Facial kinship verification is evaluated almost universally on four benchmarks:
KinFaceW-I, KinFaceW-II, TSKinFace, and Families In the Wild (FIW). We show that
the protocols in common use on these benchmarks admit two distinct forms of
optimistic bias. First, pair-level random splitting places an identity from
\emph{every} test pair into the training set: on FIW we measure
$11{,}424/11{,}424$ ($100\%$) test pairs sharing an identity with training.
Second, the standard TSKinFace pairing makes every positive pair same-family and
every negative pair cross-family, so a model can separate the classes by
recognising a shared photographic session rather than kinship; removing this cue
costs approximately $9$ ROC-AUC points. We propose a corrected protocol based on
grouped $k$-fold cross-validation in which every family is tested exactly once,
and negatives are constructed within each fold so that no family spans the
train/test boundary. Under this protocol a deliberately conventional
$0.50$M-parameter classifier attains $75.66\%$ mean accuracy and $0.8497$ mean
ROC-AUC across the four benchmarks. We additionally report a systematic negative
result: nine distinct quantum-inspired formulations, spanning similarity
measures, feature extractors, and inductive biases, each evaluated against a
capacity-matched classical control, fail to improve on their controls; four are
significantly worse. We release the protocol implementation and all evaluation
code.
\end{abstract}

\begin{IEEEkeywords}
Facial kinship verification, evaluation protocol, dataset bias, data leakage,
quantum-inspired machine learning, negative results.
\end{IEEEkeywords}

% ---------------------------------------------------------------------------
\section{Introduction}

Facial kinship verification asks whether two facial images depict biologically
related individuals. Progress is measured on a small set of public benchmarks,
and reported accuracies have risen steadily. This paper examines whether the
evaluation protocols supporting those numbers are sound.

We report three findings.

\textbf{Identity leakage.} The conventional practice of splitting a corpus at
the level of \emph{pairs} rather than \emph{identities} does not isolate the
test set. FIW contains $57{,}120$ pairs drawn from only $95$ families and $472$
distinct identities, so identities recur across pairs an average of $242$ times.
Under a random pair-level split we measure that $100\%$ of test pairs share at
least one identity with the training set. Any metric obtained this way is
optimistic by an unknown margin.

\textbf{A dataset shortcut.} In the standard TSKinFace construction, all $800$
positive pairs are drawn from within a family and all $800$ negatives across
families. Family membership therefore predicts the label perfectly, and a model
can exploit correlated capture conditions --- illumination, camera, background
--- instead of facial resemblance. We rebuild the negatives to include
same-family non-kin pairs (father versus mother, who are co-photographed but not
kin). Embedding separability drops from $0.7278$ to $0.6880$ ROC-AUC, isolating
roughly nine points attributable to the shortcut.

\textbf{A negative result for quantum-inspired metrics.} Quantum-inspired
similarity, in which embeddings are mapped to rotation angles and compared by a
simulated SWAP test, has been proposed for this task. We evaluate nine
formulations against capacity-matched classical controls. None improves on its
control; four are significantly worse. We give a structural explanation rather
than attributing the outcome to implementation.

Our contributions are: (i) quantification of both biases with reproducible
measurements; (ii) a corrected, released evaluation protocol; (iii) baseline
results under it; and (iv) a pre-registered nine-formulation negative result.

% ---------------------------------------------------------------------------
\section{Related Work}

% TODO(before submission): This section must cite the actual literature. Every
% entry below is a placeholder describing what belongs there. Do NOT populate
% it with numbers that cannot be traced to a paper you have read --- an earlier
% draft of this manuscript contained a baseline comparison table whose rows
% could not be sourced, and it has been removed.

\subsection{Kinship verification}
Metric-learning approaches, deep pairwise and graph-based models, and the
benchmarks themselves (KinFaceW, TSKinFace, FIW, and the RFIW challenge series).

\subsection{Evaluation bias in vision benchmarks}
Subject-disjoint versus random splitting in face recognition and biometrics;
shortcut learning and spurious correlation; template- and set-based protocols
such as IJB-B/C, which motivate our set-level discussion in
Section~\ref{sec:limitations}.

\subsection{Quantum-inspired machine learning}
Variational circuits, SWAP-test similarity, density-matrix classifiers, and
prior reports of quantum-inspired gains on classical data.

% ---------------------------------------------------------------------------
\section{Sources of Optimistic Bias}
\label{sec:bias}

\subsection{Identity leakage under pair-level splitting}

Let $\mathcal{P}$ be a set of labelled pairs and $\mathrm{id}(\cdot)$ the
identity depicted in an image. A split $(\mathcal{P}_{\text{tr}},
\mathcal{P}_{\text{te}})$ is \emph{identity-disjoint} if
\begin{equation}
\Big(\bigcup_{p \in \mathcal{P}_{\text{tr}}} \mathrm{id}(p)\Big) \cap
\Big(\bigcup_{p \in \mathcal{P}_{\text{te}}} \mathrm{id}(p)\Big) = \emptyset .
\end{equation}
Random splitting over $\mathcal{P}$ does not enforce this. Table~\ref{tab:leak}
reports the measured violation on FIW.

\begin{table}[t]
\caption{Identity leakage under a random $80/20$ pair-level split of FIW.}
\label{tab:leak}
\centering
\begin{tabular}{lr}
\toprule
Quantity & Value \\
\midrule
Pairs & 57{,}120 \\
Distinct families & 95 \\
Distinct identities & 472 \\
Mean appearances per identity & 242 \\
Largest family, share of all pairs & 19\% \\
\midrule
Test pairs sharing an identity with train & 11{,}424 / 11{,}424 (100\%) \\
Test pairs whose family was seen in train & 11{,}424 / 11{,}424 (100\%) \\
\bottomrule
\end{tabular}
\end{table}

\subsection{A non-obvious obstacle to correcting it}

Splitting on families is not sufficient by itself. Each \emph{negative} pair
joins two families, and treating families as vertices and negatives as edges
yields a graph over FIW with a single connected component containing all $95$
families. No family-disjoint partition of the existing pairs exists. Negatives
must therefore be discarded and regenerated \emph{within} each side after the
split. We found this empirically: a splitter that passed synthetic unit tests
failed on the real corpus for exactly this reason.

A second consequence is data loss. Because pre-existing negatives straddle any
group boundary, naively dropping them removes usable data --- on KinFaceW-I,
$173$ of $1{,}066$ pairs, leaving only $119$ test pairs. Regenerating negatives
per side instead recovers the test set to $212$ pairs on KinFaceW-I and $400$ on
KinFaceW-II.

\subsection{The TSKinFace photo-session shortcut}

Table~\ref{tab:ts} contrasts the standard TSKinFace pairing with our
shortcut-free construction, in which half the negatives are drawn within a
family as father--mother pairs: co-photographed, hence sharing every session
cue, but not kin.

\begin{table}[t]
\caption{TSKinFace pair construction. Under the standard protocol, family
membership predicts the label perfectly.}
\label{tab:ts}
\centering
\begin{tabular}{lccc}
\toprule
Construction & Same-fam.\ pos. & Same-fam.\ neg. & Cosine AUC \\
\midrule
Standard & 100\% & 0\% & 0.7278 \\
Shortcut-free (ours) & 100\% & 50\% & 0.6880 \\
\bottomrule
\end{tabular}
\end{table}

The $0.0398$ AUC reduction is measured on raw embeddings without training, so it
reflects separability available to any model, not a property of ours. We verify
separately that our trained model does not rely on a related cue: $8.2\%$ of FIW
positives come from the same source photograph (versus $0\%$ of negatives), and
excluding them changes accuracy by $+0.08$ points ($82.53\% \to 82.61\%$).

% ---------------------------------------------------------------------------
\section{Corrected Protocol}
\label{sec:protocol}

We adopt grouped $k$-fold cross-validation with $k=5$. The grouping unit is the
family, which is the unit that shares identities and capture conditions.

\begin{enumerate}
\item \textbf{Positives only are partitioned.} Families are dealt largest-first
into the currently smallest fold, which keeps folds balanced despite the highly
skewed family sizes.
\item \textbf{Negatives are generated per side}, after partitioning, from the
families present on that side.
\item \textbf{Group keys are dataset-scoped.} The family identifier
\texttt{fd\_003} occurs in both KinFaceW-I and KinFaceW-II denoting different
families; an unscoped key merges them and corrupts the disjointness check.
\item \textbf{Family contribution is capped} at $1{,}500$ pairs, preventing the
largest FIW family ($19\%$ of all pairs) from dominating whichever fold it
enters.
\item \textbf{Thresholds are calibrated on a group-disjoint validation slice} of
the training side. The test fold is scored once.
\end{enumerate}

Fold integrity is verified directly rather than assumed: zero folds exhibit
train/test group overlap, and every family is tested exactly once with no
duplicates ($95/95$ FIW, $533/533$ KinFaceW-I, $1000/1000$ KinFaceW-II,
$559/559$ TSKinFace).

\subsection{Effect on measurement stability}

Table~\ref{tab:stab} shows why a single group-disjoint split is insufficient.
Accuracy on one split varies by up to $14.2$ points with the random seed,
because FIW's test side reduces to a small number of families of very unequal
size.

\begin{table}[t]
\caption{Accuracy range across seeds (single split) versus across folds
(grouped $k$-fold).}
\label{tab:stab}
\centering
\begin{tabular}{lcc}
\toprule
Dataset & Single split & Grouped $k$-fold \\
\midrule
FIW & 61.7--75.9 (14.2) & 69.8--74.4 (4.6) \\
KinFaceW-I & 70.3--78.8 (8.5) & 72.9--79.4 (6.5) \\
KinFaceW-II & 69.0--77.8 (8.8) & 71.0--77.0 (6.0) \\
TSKinFace & 75.1--84.1 (9.0) & 75.7--81.9 (6.2) \\
\midrule
Overall spread & 22.5 & 12.1 \\
\bottomrule
\end{tabular}
\end{table}

% ---------------------------------------------------------------------------
\section{Reference Model}
\label{sec:model}

To measure the protocol rather than an architecture, we use a deliberately
conventional model. Frozen FaceNet embeddings ($512$-d) pass through a shared
Siamese encoder; kinship is symmetric, so only symmetric pair features are
formed --- sum, absolute difference, elementwise product, and cosine --- and
concatenated with a one-hot relation code before a three-layer MLP. Predictions
are averaged over both argument orders, making the model exactly symmetric. The
model has $497{,}881$ parameters.

Its plainness is deliberate: findings in Sections~\ref{sec:bias} and
\ref{sec:protocol} cannot then be attributed to architectural novelty. We note
that a linear probe on the same frozen embeddings reaches $0.702$ ROC-AUC
against this model's $0.810$ on FIW, confirming the head performs non-trivial
work and that the frozen backbone is not the binding constraint.

% ---------------------------------------------------------------------------
\section{Results}
\label{sec:results}

\begin{table}[t]
\caption{Grouped 5-fold results. Every family tested exactly once.}
\label{tab:main}
\centering
\begin{tabular}{lcc}
\toprule
Dataset & Accuracy (95\% CI) & ROC-AUC \\
\midrule
KinFaceW-I & $76.92 \pm 2.19$ & $0.8752 \pm 0.0146$ \\
KinFaceW-II & $74.25 \pm 1.95$ & $0.8320 \pm 0.0190$ \\
FIW & $72.21 \pm 1.55$ & $0.8101 \pm 0.0175$ \\
TSKinFace (shortcut-free) & $79.25 \pm 2.32$ & $0.8814 \pm 0.0192$ \\
\midrule
\textbf{Mean} & $\mathbf{75.66}$ & $\mathbf{0.8497}$ \\
\bottomrule
\end{tabular}
\end{table}

Table~\ref{tab:main} reports results under the corrected protocol. We
deliberately omit a comparison against published accuracies. Those numbers are
obtained under each dataset's official protocol, which permits the leakage and
shortcut quantified in Section~\ref{sec:bias}; presenting a ranking against them
would mislead in both directions. We regard establishing comparable baselines
under the corrected protocol as work the community must do collectively, and we
release our implementation to that end.

\subsection{Per-relation analysis}

Performance is markedly uneven across relation types. On FIW, father--daughter
reaches $0.9933$ ROC-AUC while mother--son reaches $0.7124$ --- and mother--son
is the largest subset ($5{,}296$ pairs). Averaged over datasets the ordering is
mother--son ($0.8288$), mother--daughter ($0.8973$), father--son ($0.9172$),
father--daughter ($0.9312$). We verified that this is not a labelling artefact:
gender fields in the FIW metadata are complete and well-formed across the
families we audited. Cross-gender parent--child verification appears
intrinsically harder, and we regard it as the most promising target for future
architectural work.

% ---------------------------------------------------------------------------
\section{Set-Level and Triadic Representations}
\label{sec:setlevel}

Sections~\ref{sec:bias}--\ref{sec:results} concern how the task is
\emph{measured}. We now examine two ways in which the task is conventionally
\emph{posed} that discard available evidence.

\subsection{Identity photo sets}

The standard protocol scores one photograph against one photograph. FIW,
however, supplies a median of five photographs per identity, so the
conventional formulation discards most of the evidence for each person. We
represent each identity by its photo set and form symmetric set statistics: the
$\ell_2$-normalised centroid, the within-set spread ($1$ minus mean pairwise
cosine), the set size, and the maximum, minimum, mean and standard deviation of
the cross-set similarity matrix.

A set of size one must reduce exactly to the single-image case: the centroid is
the embedding, the spread is zero, and every cross-set statistic collapses to
the single cosine. We assert this in the test suite rather than assume it,
which matters because the three remaining benchmarks store exactly one
photograph per person.

\begin{table}[t]
\caption{Set-level representation under grouped 5-fold. Gains appear only where
photo sets exist; single-photograph corpora are bit-identical to the baseline,
as required.}
\label{tab:setlevel}
\centering
\begin{tabular}{lcccc}
\toprule
Dataset & Set size & Single & Set-level & $\Delta$ \\
\midrule
KinFaceW-I & 1.00 & 0.7210 & 0.7210 & $0.0000$ \\
KinFaceW-II & 1.00 & 0.7672 & 0.7672 & $0.0000$ \\
TSKinFace & 1.00 & 0.7306 & 0.7306 & $0.0000$ \\
FIW & 7.30 & 0.7296 & \textbf{0.8066} & $\mathbf{+0.0771}$ \\
\bottomrule
\end{tabular}
\end{table}

Table~\ref{tab:setlevel} reports the outcome. On FIW the representation is worth
$+0.0771$ ROC-AUC, closely matching an untrained embedding-level estimate of
$+0.071$, which indicates the gain is intrinsic to the representation rather
than a product of additional model capacity. The three single-photograph corpora
are unchanged to machine precision.

We stress the narrowness of this result: it is a large gain on one benchmark and
exactly zero on three, and any aggregate over the four would obscure that. The
practical implication is that kinship verification deployments in which several
photographs of a person are available --- the common case --- are leaving a
substantial margin unclaimed, and that benchmarks storing a single photograph
per identity cannot measure this.

\subsection{Triadic structure}

TSKinFace supplies father, mother and child jointly, but the pairwise protocol
decomposes each triad into two independent comparisons. We instead score the
triad: alongside $\cos(F,C)$ and $\cos(M,C)$ we compute the best convex
parental mixture $\max_\alpha \cos\big(\widehat{\alpha F + (1-\alpha)M},\,
C\big)$, the maximising $\alpha$, the gain of that mixture over the better
single parent, and the parental similarity $\cos(F,M)$.

The last quantity is not expressible in a pairwise formulation and is the
strongest single contributor: a child resembling one parent is weaker evidence
of kinship when the two parents already resemble each other.

\begin{table}[t]
\caption{Triadic scoring on TSKinFace, grouped 5-fold, folds assigned by the
parents' family.}
\label{tab:triadic}
\centering
\begin{tabular}{lcc}
\toprule
Arm & Accuracy (95\% CI) & ROC-AUC \\
\midrule
Best single parent & $67.62 \pm 2.69$ & $0.7569$ \\
Triadic & $\mathbf{69.05 \pm 2.65}$ & $\mathbf{0.7893}$ \\
Triadic $+$ phase sweep & $69.50 \pm 2.41$ & $0.7887$ \\
\bottomrule
\end{tabular}
\end{table}

Triadic scoring improves ROC-AUC by $+0.0324$ over the best single parent
(paired $t$-test over folds, $t = 14.93$, $p = 0.0001$; the gain is positive in
all five folds). Negatives retain the true parents and substitute a child from
another family, so the label depends on kinship rather than on any property of
the parent pair.

% ---------------------------------------------------------------------------
\section{A Negative Result for Quantum-Inspired Metrics}
\label{sec:quantum}

\subsection{Formulations tested}

We evaluated nine formulations, chosen to span the three roles a quantum
component can occupy: a similarity measure, a feature extractor, and an
inductive bias. Each is compared against a control matched for capacity or
intent (Table~\ref{tab:quantum}).

\begin{table}[t]
\caption{Quantum-inspired formulations against matched classical controls.}
\label{tab:quantum}
\centering
\small
\begin{tabular}{llc}
\toprule
Formulation & Control & Outcome \\
\midrule
SWAP fidelity (5 seeds) & classical head & $p=0.945$ \\
SWAP fidelity (20 folds) & classical head & $p=0.982$ \\
Density fidelity $\mathrm{Tr}(\rho_1\rho_2)$ & mean-pooling & $-8.7$ AUC \\
Interference phase sweep & classical mixture & $+0.1$ AUC \\
Entanglement entropy & cosine & $0.514$ alone \\
POVM head & capacity-matched head & $-5.5$ AUC \\
Purity regulariser & decorrelation penalty & $-0.9$, $p=0.047$ \\
\bottomrule
\end{tabular}
\end{table}

Under grouped $k$-fold with $20$ paired observations, enabling the SWAP-test
branch yields $75.66\%$ against $76.34\%$ with it disabled (accuracy
$p = 0.158$; ROC-AUC $p = 0.982$). The branch also costs $6$--$14\times$ runtime.

\subsection{A cautionary observation}

The purity regulariser initially measured $+1.3$ ROC-AUC on a single seed, with
reduced variance --- a result consistent with effective regularisation. Repeated
across nine paired runs the sign reversed, and the method proved significantly
\emph{worse} than no regulariser ($-0.0089$, $p = 0.047$). We report this
because the single-seed reading was of the kind that is readily published, and
it illustrates why capacity-matched controls and repeated trials are necessary
when evaluating claims of this type.

\subsection{Two further formulations}

The set-level and triadic representations each admit a natural quantum reading,
and we tested both as pre-registered controls.

A photo set is a mixed state $\rho = \frac{1}{N}\sum_i
|\psi_i\rangle\langle\psi_i|$, and $\mathrm{Tr}(\rho_1 \rho_2)$ is the
corresponding fidelity. Added to the set features it changes ROC-AUC by
$-0.0006$, $-0.0005$, $+0.0001$ and $-0.0006$ on the four benchmarks.

A child explained by two parents is naturally written as a coherent
superposition with a relative phase, of which the classical convex mixture is
the zero-phase special case. Sweeping the phase changes ROC-AUC by $-0.0006$
($p = 0.147$).

This brings the total to nine formulations, none of which improves on its
classical control.

\subsection{Why the approach fails}

The failures are structural rather than implementational. Interference-based
formulations reduce to their classical counterparts because $\ell_2$-normalised
FaceNet embeddings carry no phase structure to exploit; the phase sweep in
Table~\ref{tab:quantum} is maximised at the classical mixture. Density-matrix
formulations require estimating a $d \times d$ operator from the images
available for one identity --- a median of five --- so the resulting $\rho$ is
dominated by its leading eigenvector, which is the mean, while accumulating
noise from tail directions that are not estimable. This is why
$\mathrm{Tr}(\rho_1\rho_2)$ scores \emph{below} simple mean-pooling.

\subsection{An efficiency contribution}

Independently of the above, we contribute an exact reformulation of the
simulated SWAP test. The reference implementation materialises a $2^n$
statevector and permutes all $n$ axes once per qubit, incurring kernel-launch
overhead that leaves the accelerator idle. Applying rotations by reshaping to
$(B, \text{left}, 2, \text{right})$ and precomputing the CNOT chain and $R_z$
layer as a single diagonal phase vector yields a $380\times$ speedup on GPU
($192 \to 73{,}094$ pairs/s) with identical output and gradients, verified by
test. We note that a closed-form product over qubits is \emph{not} available
here: the shared CNOT chain correlates the $R_z$ phases, so the overlap does not
factorise. The gain arises from removing overhead, not from approximation.

This is a quantum-versus-quantum speedup. Simulating a circuit is necessarily
more costly than the classical operation it parallels, and we measure the
quantum branch at $5$--$14\times$ slower than the classical path at every qubit
count from $2$ to $12$.

% ---------------------------------------------------------------------------
\section{Limitations}
\label{sec:limitations}

\textbf{Small test sets.} KinFaceW-I and KinFaceW-II yield $212$ and $400$ test
pairs per fold, giving confidence intervals of roughly $\pm 2$ points. FIW and
TSKinFace are firmer.

\textbf{Single backbone.} All results use frozen FaceNet embeddings. Whether the
biases we quantify interact with backbone choice is untested.

\textbf{Set-level gains are confined to one benchmark.} Only FIW stores
multiple photographs per identity, so the $+0.0771$ ROC-AUC of
Section~\ref{sec:setlevel} is measured on a single corpus; the other three are
unchanged by construction. Whether the effect holds at other set sizes, and on
corpora we could not test, is open.

\textbf{Negative results are bounded.} We tested nine formulations; we do not
claim no quantum-inspired method can help on this task, only that these nine,
covering the three natural roles, do not.

% ---------------------------------------------------------------------------
\section{Conclusion}

Two properties of standard facial kinship verification protocols inflate
reported performance: pair-level splitting leaks identities into the test set at
a measured rate of $100\%$ on FIW, and the standard TSKinFace pairing admits a
photo-session shortcut worth approximately nine ROC-AUC points. We provide a
corrected grouped $k$-fold protocol, verify its integrity directly, and report
baselines under it. We further report that nine quantum-inspired formulations,
each against a matched control, fail to improve on classical baselines, and give
a structural account of why. All code is released so that the protocol can be
adopted and the results checked.

% ---------------------------------------------------------------------------
\section*{Reproducibility}

Every numeric claim derives from JSON artefacts in \texttt{results/honest/}.
The protocol is implemented in \texttt{src/splits.py}, \texttt{src/kfold.py},
and \texttt{src/ts\_pairs.py}; results are regenerated by
\texttt{scripts/evaluation/run\_kfold.py}. The implementation is covered by 61
tests, including assertions that no group spans a fold boundary and that the
fast SWAP-test reformulation matches the reference simulator in both value and
gradient.

% ---------------------------------------------------------------------------
\begin{thebibliography}{1}
% TODO(before submission): populate with references actually consulted.
% An earlier draft cited baselines that could not be traced to a source; those
% rows were removed rather than reproduced. Add entries only for papers read.
\bibitem{placeholder}
Reference list to be completed.
\end{thebibliography}

\end{document}
